\section{Do you need this yet?}
\label{sec:ch01-do-you-need-this}

I have spent five sections describing a disease and naming companies that
caught it. Here is the part that gets left out of conference talks: plenty of
organisations operating at genuinely large scale looked at federation, and
decided against it, and were right.

\index{federation!when not to adopt}%
The most instructive example is Twitter's Core API Platform team, writing in
2021 about a schema that, by their own description, defined \enquote{the
entirety of our api data model} with roughly a thousand object types, a hundred
input types, and three hundred mutation fields, growing daily as \enquote{several
hundred developers across the company add fields and types}. That graph served
around 500,000 requests per second, and one of their most frequent queries was
2,500 lines long and returned 50,000 fields. When someone asked them directly
why they had not federated, the answer was that they did not need to:
\enquote{our underlying data layer (known as strato) is itself federated}, which
gave them \enquote{most of the perks of federation while also allowing us to
have a unified schema with strong network effects and without too many
performance gotchas}~\autocite{twitter2021unified}.

Read that last clause again, because it is a platform engineer at half a million
requests per second saying that a federated graph has performance
characteristics his unified schema does not have to deal with. It is also a
reminder that federation of the graph is not the only place you can put a seam.
They put theirs one layer down, in the data access tier, which is one of the
alternatives chapter~\ref{ch:do-you-actually-need-federation} takes seriously.

Zalando made a similar call for similar reasons. They run a unified graph across
their domains, but as a single service rather than a composed one, and Aditya
Pratap Singh was explicit that this is a trade rather than an oversight: they
\enquote{gain by keeping a single Graph in terms of tooling, deployment and
governance}~\autocite{singh2021zalando}. Tooling, deployment, and governance are
exactly the three things federation makes harder, and they are not glamorous
enough to make it into most architecture decisions.

\index{GitHub!GraphQL API}%
Then there is GitHub, which is the example I reach for when someone argues that
a schema can simply become too large for one service. I counted the published
GitHub GraphQL schema while writing this chapter: 1,807 named types and 6,809
fields on object and interface types, with 32 root query fields and 274
mutations~\autocite{github2026schema}. It is not versioned. It carries 1,258
deprecations, most of them with a removal date attached, governed by a published
policy that promises at least three months' notice and lands breaking changes on
quarter boundaries~\autocite{github2026breaking}. That is one schema, an order
of magnitude larger than most schemas anybody reading this owns, kept coherent
by a deprecation process rather than a distributed architecture. GraphQL's own
documentation notes that Meta, where the language was invented, has run a
monolithic GraphQL API since 2012~\autocite{graphqlorg2026federation}.

Schema size, then, is not the trigger. Neither is request volume. The trigger is
the number of teams that need to change the same part of the schema in the same
week, which is why the exercise at the end of this chapter asks you to count
that rather than to count fields.

\subsection*{What it costs}

If you do federate, here is the bill. I would rather you see it now than
discover it in Part~VI.

The router and the schema registry become production dependencies with their own
failure modes. Booking.com published the most useful account of this I know:
running close to a thousand gateway instances that fetched their schema roughly
8.3 million times a day, they hit a 24-hour window in which the failure rate for
schema delivery reached about 2.4 per cent against an expected 0.05 per cent,
and a gateway that cannot fetch a schema cannot roll
out~\autocite{ernst2022booking}. None of that is a criticism of federation. It
is what happens when composition becomes something your deploys depend on, and
it is work you were not doing before.

The performance surface is larger and harder to see into. At Netflix, where two
subgraph services between them handle over a million GraphQL queries per second,
Stephen Chambers described fixing a defect in the federation query-planning
layer, and the fix cut requests per second by 20 per cent and saved hundreds of
thousands of dollars~\autocite{chambers2026netflix}. That is a cost which had
been running silently until somebody went looking for it, in a layer that did
not exist before they federated.

The expertise cost is distributed rather than eliminated, which is easy to
misread as a saving. Netflix noted that in the monolith era only the Studio API
team needed to know GraphQL deeply, whereas afterwards \enquote{every DGS team
needs to build expertise in GraphQL}~\autocite{shikhare2020netflix2}. Airbnb,
who built a different answer to the same problem, put the structural version of
this well: \enquote{Federation distributes development by distributing
servers}~\autocite{tanner2026viaduct}. Hundreds of contributing teams means
hundreds of services to run, monitor, and page someone about.

\index{federation!maturity}%
It is worth noting who else says this. Netflix, the reference customer for the
entire idea, closed their own account of it by warning that federation
\enquote{is early in its maturity lifecycle and may not be the best fit for
every team or organization}, and that learning it, running a federation layer,
and carrying out the migration together \enquote{requires high commitment from
partner teams and seamless cross-functional
collaboration}~\autocite{shikhare2020netflix2}. GraphQL's own documentation is
blunter, saying federation \enquote{requires substantial infrastructure support,
including a dedicated team to manage the gateway, schema registry} and that
therefore \enquote{it's crucial to consider whether your organization truly
needs this level of complexity}~\autocite{graphqlorg2026federation}. When the
language's own home page tells you to think harder, that is not marketing.

\subsection*{A first pass}

So, a rough test, pending the real one in
chapter~\ref{ch:do-you-actually-need-federation}. Federation starts to earn its
cost when several teams own genuinely different domains inside one schema, when
those teams collide on the same types often enough that you can remember
specific incidents, when they need to deploy on their own schedules, and when
you already have the operational maturity to run another tier in production
without flinching. If instead you have one team, or several teams that rarely
touch each other's types, or nobody on call for a gateway, then the honest
answer is that your problems have cheaper solutions: a style guide, a
\code{CODEOWNERS} file, a deprecation policy like GitHub's, a modular monolith
with real internal boundaries.

Martin Fowler's advice to build the monolith first and split it once you can see
where the seams belong was published in 2015 and has aged
well~\autocite{fowler2015monolithfirst}. It is not unanimous, and the dissent is
worth reading: Stefan Tilkov argues on the same site that systems intended from
the start to be distributed are hard to retrofit into that shape
later~\autocite{tilkov2015dontstart}. For graphs specifically I still lean
Fowler's way, because a schema is unusually easy to carve later. The entity keys
that federation needs are mostly the identifiers you already have, and
chapter~\ref{ch:strangling-the-monolith} does exactly this carving without
downtime.

One last observation, offered as a caution rather than a conclusion. While
researching this chapter I looked for companies that had adopted federation and
then removed it, and found essentially none. That silence is real, but its cause
is not knowable from the outside. It might mean federation works. It might mean
that once a graph is composed, unpicking it costs more than living with it,
which is an argument for deciding carefully rather than a reassurance. Either
reading points the same direction: this is a decision worth making deliberately,
in advance, with your own numbers.
